\documentclass[10pt]{article}
\usepackage[a4paper,landscape,top=12mm,bottom=12mm,left=13mm,right=13mm]{geometry}
\usepackage[T1]{fontenc}
\usepackage[utf8]{inputenc}
\usepackage[portuguese]{babel}
\usepackage{graphicx,xcolor,amsmath,microtype,booktabs}
\usepackage[hidelinks]{hyperref}
\pagestyle{empty}
\setlength{\parindent}{0pt}
\setlength{\parskip}{4pt}

\definecolor{tinta}{HTML}{272727}
\definecolor{cinza}{HTML}{4D4D4D}
\definecolor{claro}{HTML}{767676}
\definecolor{vermelho}{HTML}{B64342}
\definecolor{azul}{HTML}{0F4D92}
\definecolor{teal}{HTML}{42949E}
\color{tinta}

% cabeçalho de página de análise
\newcommand{\cab}[3]{%
  {\color{claro}\footnotesize\textsf{#1}}\\[-2pt]
  {\LARGE\textbf{#2}}\\[2pt]
  {\color{cinza}\small #3}\par\vspace{4pt}
  {\color{claro}\rule{\linewidth}{0.4pt}}\par\vspace{5pt}}

% rótulo de bloco no painel de texto
\newcommand{\bloco}[1]{\vspace{2pt}{\color{vermelho}\footnotesize\textbf{#1}}\par\vspace{-2pt}}
\newcommand{\ressalva}[1]{{\color{claro}\scriptsize #1}\par\vspace{1pt}}
\newcommand{\lit}[1]{\vspace{2pt}{\color{azul}\footnotesize\textbf{Na literatura}}\par\vspace{-2pt}{\color{cinza}\scriptsize #1}\par\vspace{1pt}}

% página: figura à esquerda, leitura à direita
\newcommand{\analise}[6]{%
  \cab{#1}{#2}{#3}
  \begin{minipage}[t]{0.588\linewidth}\vspace{0pt}
    \includegraphics[width=\linewidth]{#4}\par\vspace{6pt}
    {\color{claro}\rule{\linewidth}{0.4pt}}\par\vspace{4pt}
    {\scriptsize\setlength{\parskip}{0pt}\color{cinza}#6\par}
  \end{minipage}\hfill
  \begin{minipage}[t]{0.383\linewidth}\vspace{0pt}%
    \footnotesize\setlength{\parskip}{3pt}#5
  \end{minipage}
  \clearpage}

\newcommand{\analisec}[6]{%
  \cab{#1}{#2}{#3}
  \begin{minipage}[t]{0.588\linewidth}\vspace{0pt}
    \includegraphics[width=\linewidth]{#4}\par\vspace{6pt}
    {\color{claro}\rule{\linewidth}{0.4pt}}\par\vspace{4pt}
    {\scriptsize\setlength{\parskip}{0pt}\color{cinza}#6\par}
  \end{minipage}\hfill
  \begin{minipage}[t]{0.383\linewidth}\vspace{0pt}%
    \scriptsize\setlength{\parskip}{2.5pt}#5
  \end{minipage}
  \clearpage}


\newcommand{\analisea}[6]{%
  \cab{#1}{#2}{#3}
  \begin{minipage}[t]{0.515\linewidth}\vspace{0pt}
    \includegraphics[width=\linewidth]{#4}\par\vspace{6pt}
    {\color{claro}\rule{\linewidth}{0.4pt}}\par\vspace{4pt}
    {\scriptsize\setlength{\parskip}{0pt}\color{cinza}#6\par}
  \end{minipage}\hfill
  \begin{minipage}[t]{0.455\linewidth}\vspace{0pt}%
    \scriptsize\setlength{\parskip}{2.5pt}#5
  \end{minipage}
  \clearpage}



\newcommand{\legI}{\textbf{Fig. 1 | O corte de geração fotovoltaica no SIN é sistêmico, de razão energética e incide sobre toda a janela solar.} \textbf{a}, Perfil diurno médio da potência fotovoltaica das usinas com apuração de \textit{constrained-off} no Sistema Interligado Nacional. A área azul é a geração verificada e a área vermelha é a energia cortada, definida como \texttt{max(referência $-$ verificada, 0)} nos intervalos com restrição apurada; a soma define a potência que estaria disponível. Médias sobre todos os intervalos de 30 min do período; horário de Brasília, não corrigido para hora solar local. \textbf{b}, Taxa mensal de corte por subsistema, definida como energia cortada dividida pela energia disponível (verificada + cortada). \textbf{c}, Energia total cortada por razão de restrição, decomposta em origem sistêmica e local; rótulos indicam o total por razão. \textbf{d}, Distribuição espacial das usinas e conjuntos fotovoltaicos com apuração de \textit{constrained-off}; a cor codifica a taxa de corte acumulada e o tamanho, a capacidade instalada. A escala de cor satura no percentil 97 para preservar contraste. Contornos estaduais: IBGE. Dados: ONS, dados abertos, restrição por \textit{constrained-off} de usinas fotovoltaicas e fator de capacidade. Dados-fonte em \texttt{data/processed/fig1\_source\_*.csv}.}
\newcommand{\legII}{\textbf{Fig. 2 | A MMGD brasileira é majoritariamente sub-pixel no Sentinel-2, e a série de conexões carrega uma descontinuidade regulatória.} \textbf{a}, Distribuição acumulada da área do arranjo das instalações fotovoltaicas de micro e minigeração distribuída registradas na ANEEL. Linhas tracejadas marcam a área de um pixel de cada sensor; o rótulo indica a fração de instalações menor que um pixel. \textbf{b}, Fração de instalações que atinge o piso operacional de 4 pixels, por classe de consumo e sensor. \textbf{c}, Conexões mensais; a linha tracejada marca o fim da janela de transição da Lei 14.300/2022 e o marcador, o mês de pico. \textbf{d}, Participação da classe rural no total de instalações de cada unidade da federação. Contornos: IBGE. Dados: ANEEL, relação de empreendimentos de geração distribuída e informações técnicas fotovoltaicas. Dados-fonte em \texttt{data/processed/fig2\_source\_*.csv}.}
\newcommand{\legIII}{\textbf{Fig. 3 | Dois ativos solares em Minas Gerais diferem 3,0$\times$ na taxa de corte, e a diferença não se reproduz como padrão no parque.} \textbf{a}, Perfil diurno médio de potência, em pu da capacidade instalada, na janela comum aos dois ativos. Azul, geração verificada; vermelho, energia cortada. \textbf{b}, Custo Marginal de Operação do subsistema Sudeste ponderado pela energia, durante os intervalos de corte e durante a geração. \textbf{c}, Taxa de corte mensal dos dois ativos. \textbf{d}, Pixels Sentinel-2 ocupados pela área de módulos, comparados à instalação mediana de MMGD; escala logarítmica, área de módulos estimada a 5,12 m\textsuperscript{2}/kWp. \textbf{e}, Taxa de corte contra capacidade renovável variável num raio de 100 km, para os conjuntos com apuração completa na janela de referência da frota; os dois alvos em destaque. Dados: ONS (constrained-off, fator de capacidade, CMO, relacionamento conjunto-usina) e ANEEL (SIGA). Dados-fonte em \texttt{data/processed/fig3\_source\_*.csv}.}
\newcommand{\legIV}{\textbf{Fig. 4 | A queda do fator de capacidade de Sol do Cerrado acompanha a energização do entorno; o controle não mostra o mesmo.} \textbf{a}, Jaíba (MG). Área cinza, capacidade fotovoltaica acumulada em operação num raio de 100 km (eixo esquerdo); linha, fator de capacidade mensal verificado de Sol do Cerrado (eixo direito); faixa vermelha, parcela cortada apurada pelo ONS, disponível apenas a partir de abril de 2024 (linha tracejada). \textbf{b}, São Gonçalo do Abaeté (MG), com a mesma construção para Três Marias 3. \textbf{c}, Fator de capacidade mensal de Sol do Cerrado em 2023 e 2025, mês contra mês, para separar restrição de sazonalidade. \textbf{d}, Potência acumulada de micro e minigeração distribuída nos dois municípios; os 20,4 MW de Jaíba equivalem a 3,0\% da capacidade de Sol do Cerrado. Dados: ONS (fator de capacidade, constrained-off) e ANEEL (SIGA, MMGD). Dados-fonte em \texttt{data/processed/fig4\_source\_*.csv}.}
\newcommand{\legV}{\textbf{Fig. 5 | Normalizada pela irradiância medida sobre a usina, a perda de Sol do Cerrado permanece; o recurso solar não a explica.} \textbf{a}, Razão de desempenho mensal, definida como energia gerada por kW instalado dividida pela irradiância global horizontal diária; tracejado horizontal marca a média de cada ano. \textbf{b}, Irradiância global horizontal mensal sobre cada usina (linha cheia) e a mesma sob céu claro (tracejado). \textbf{c}, Decomposição da queda do fator de capacidade de Sol do Cerrado entre 2023 e 2025: observado em 2023, esperado em 2025 caso apenas a irradiância tivesse mudado, e observado em 2025. \textbf{d}, Fator de capacidade mensal contra irradiância mensal em Sol do Cerrado, por ano. Dados climáticos: NASA POWER (MERRA-2 + SYN1DEG), série diária no ponto de cada usina. Dados de energia: ONS. Dados-fonte em \texttt{data/processed/fig5\_source\_*.csv}.}
\newcommand{\legVI}{\textbf{Fig. 6 | A geração de referência do ONS é bem calibrada; o viés aparente da Fig. 5 vinha do contrafactual, não do operador.} \textbf{a}, Distribuição do viés mediano por usina, definido como a razão entre o potencial implicado pelo ONS (verificado + cortado) e o potencial previsto pela irradiância local usando a razão de desempenho calibrada nos dias sem restrição da própria usina; tracejado em 1 marca ausência de viés. \textbf{b}, Viés contra a intensidade mediana de corte nos dias restritos. \textbf{c}, Validação da calibração: fator de capacidade contra irradiância nos dias sem restrição de Sol do Cerrado, com a reta ajustada pela razão de desempenho limpa. \textbf{d}, Razão de desempenho anual de Sol do Cerrado contra o desempenho limpo calibrado; o rótulo interno indica o corte implícito de cada ano. Dados: ONS (fator de capacidade, constrained-off) e NASA POWER. Dados-fonte em \texttt{data/processed/fig6\_source\_*.csv}.}
\newcommand{\legVII}{\textbf{Fig. 7 | A construção de usinas fotovoltaicas é detectável e datável por Sentinel-2; a operação não é.} \textbf{a}, Razão entre o brilho no vermelho dentro do arranjo e num anel de controle de 1,5 a 3 km, por usina (linhas cinza) e mediana entre usinas (linha vermelha, faixa entre quartis), alinhadas pela data de entrada em operação do SIGA. \textbf{b}, Sentinel-2 em cor natural sobre o complexo Sol do Cerrado, em escala de refletância fixa entre as datas; círculos marcam as 17 usinas do SIGA. \textbf{c}, Comparação pareada da razão de brilho entre fases, uma linha por usina, com a mediana em destaque; \texttt{p} de Wilcoxon pareado. \textbf{d}, Diferença entre o ano de pico do brilho relativo e o ano de entrada em operação; barras em vermelho indicam erro de até um ano. Imagens: Sentinel-2 L2A (ESA), via catálogo STAC público. Registro: ANEEL/SIGA. Dados-fonte em \texttt{data/processed/fig7\_source\_*.csv}.}
\newcommand{\legVIII}{\textbf{Fig. 8 | O corte fotovoltaico por excesso de oferta coincide com térmica inflexível, mas apenas em carga baixa.} \textbf{a}, Perfil diurno médio nacional. Área azul, geração fotovoltaica verificada; área vermelha empilhada, energia cortada por razão energética (\texttt{ENE}); linhas, geração térmica separada em parcela econômica (por ordem de mérito, quando o CVU da usina é menor que o CMO) e parcela fora do mérito (geração total menos ordem de mérito). A anotação registra a carga líquida --- carga do SIN menos geração fotovoltaica --- na madrugada e no pico solar. \textbf{b}, Composição da geração térmica fora do mérito nas horas com e sem corte por razão energética, decomposta pelos motivos não-econômicos disjuntos do dicionário do ONS; rótulo indica o total. \textbf{c}, Correlação de Spearman entre corte por razão energética e geração térmica fora do mérito, estratificada por quintil de carga do SIN, restrita às horas de 9 h a 15 h; barras cinza indicam associação não significativa; abaixo do eixo, o corte médio de cada quintil. \textbf{d}, Controle negativo: taxa diária de corte por razão energética contra a energia armazenada do SIN, com a mediana por quartil de armazenamento. Dados: ONS, dados abertos --- restrição por \textit{constrained-off} fotovoltaico, geração térmica por motivo de despacho, curva de carga e energia armazenada por subsistema. Dados-fonte em \texttt{data/processed/fig8\_source\_*.csv}.}
\newcommand{\legIX}{\textbf{Fig. 9 | O piso térmico inflexível está onde o corte não está, e a coincidência nacional da Fig. 8 não sustenta leitura causal local.} \textbf{a}, Potência média nas horas solares (9 h--15 h) por subsistema: corte fotovoltaico por razão energética, geração térmica fora da ordem de mérito e a parcela desta que é inflexibilidade pura; abaixo do eixo, a razão entre o corte e o piso inflexível local. \textbf{b}, Correlação de Spearman entre corte por razão energética e térmica fora do mérito \textbf{do mesmo subsistema}, por quintil de carga do próprio subsistema; tracejado no limiar de magnitude declarado. \textbf{c}, A mesma correlação medida contra a térmica do próprio subsistema e contra a do outro subsistema com corte material. \textbf{d}, Distribuição acumulada do fluxo Nordeste$\rightarrow$Sudeste nas horas de corte alto (quartil superior do corte do Nordeste) e nas demais; tracejado no percentil 95 do fluxo, usado como referência empírica de teto por não haver limite operativo publicado. Dados: ONS, dados abertos --- \textit{constrained-off} fotovoltaico, geração térmica por motivo de despacho, curva de carga e intercâmbio nacional. Dados-fonte em \texttt{data/processed/fig9\_source\_*.csv}.}
\newcommand{\legX}{\textbf{Fig. 10 | A distância de rede até a demanda explica parte da alocação do corte que covariáveis locais não explicavam.} \textbf{a}, Grafo de transmissão em operação --- subestações e linhas da EPE, separadas em 500 kV ou mais e abaixo disso --- com as usinas fotovoltaicas da amostra posicionadas no seu ponto de conexão; a cor codifica a taxa de corte na janela de referência da frota e o tamanho, a capacidade instalada. Estrelas marcam as nove regiões metropolitanas usadas como centros de carga; a moldura cobre a região que contém a amostra, e 353 das 1.896 linhas do grafo ficam fora dela e não são desenhadas --- o recorte é de desenho, e todas as arestas participam do cálculo dos caminhos mínimos. \textbf{b}, Correlação de Spearman entre a taxa de corte e cada uma das oito covariáveis de rede pré-declaradas; em vermelho as que sobrevivem à correção de Holm para oito hipóteses. \textbf{c}, Taxa de corte contra a distância de rede ao centro de carga mais próximo, por subsistema, com a mediana por quartil de distância. \textbf{d}, As duas covariáveis sobreviventes medidas \textbf{dentro} de cada subsistema, para separar efeito de rede de efeito de subsistema. Dados: EPE, camadas de linhas de transmissão e subestações em operação; ONS, restrição por \textit{constrained-off} fotovoltaico e fator de capacidade. Dados-fonte em \texttt{data/processed/fig10\_source\_*.csv}.}
\newcommand{\legXI}{\textbf{Fig. 11 | O achado da Fig. 10 não se replica em eólica, e a dissociação espelha a origem do corte de cada fonte.} \textbf{a}, Correlação de Spearman entre a taxa de corte e as oito covariáveis de rede pré-declaradas na Fig. 10, medidas separadamente em fotovoltaica e eólica; a cor cheia marca as que sobrevivem à correção de Holm aplicada dentro de cada fonte, e o cinza as que não sobrevivem. \textbf{b}, Energia cortada decomposta por razão e origem da restrição, conforme classificação operativa do ONS, para cada fonte. \textbf{c}, As duas covariáveis-chave medidas apenas no Nordeste, onde as duas amostras coexistem, para separar tecnologia de geografia. \textbf{d}, Taxa de corte eólica contra o número de conjuntos no mesmo ponto de conexão, com a mediana por valor. Dados: ONS, restrição por \textit{constrained-off} eólico e fotovoltaico e fator de capacidade; EPE, linhas de transmissão e subestações em operação. Dados-fonte em \texttt{data/processed/fig11\_source\_*.csv}.}
\newcommand{\legXII}{\textbf{Fig. 12 | Dos quatro achados de rede da série, um passa nos três testes de robustez e o principal da Fig. 10 passa em apenas um.} \textbf{a}, Distribuição bootstrap do coeficiente de Spearman de cada achado, com 4.000 reamostras com reposição; a barra marca o intervalo de 95\% e o ponto, o valor da amostra. \textbf{b}, Coeficiente calculado separadamente nas duas metades de uma partição estrutural --- usinas cuja barra de conexão consta do cadastro de rede de 230 kV ou mais do ONS, e as demais --- com o valor $p$ do teste $z$ de Fisher para a diferença entre metades. \textbf{c}, O mesmo coeficiente calculado no grafo da EPE, usado nas Figs. 10 e 11, e no grafo do ONS, construído por número de barra; abaixo de cada barra, o tamanho da amostra no grafo alternativo, ou a marca de que o critério não se aplica --- a capacidade da usina não deriva do grafo e não tem versão alternativa. \textbf{d}, Placar: quantos dos critérios aplicáveis cada achado passa; o traço marca o critério não aplicável. Dados: Figs. 10 e 11; ONS, linhas de transmissão e subestações da Rede de Operação e fator de capacidade. Dados-fonte em \texttt{data/processed/fig12\_source\_robustez.csv}.}

\begin{document}

%================================================================ CAPA
\vspace*{6mm}
{\color{claro}\textsf{ANÁLISE DE DADOS ABERTOS DO SETOR ELÉTRICO BRASILEIRO}}\\[3mm]
{\fontsize{30}{34}\selectfont\textbf{O corte fotovoltaico no SIN}}\\[2mm]
{\fontsize{16}{20}\selectfont\color{cinza} Doze análises, cinco delas voltadas sobre as anteriores, e um achado que não resistiu ao próprio teste}\\[6mm]
{\color{claro}\rule{\linewidth}{0.4pt}}\\[4mm]

\begin{minipage}[t]{0.47\linewidth}\vspace{0pt}
\bloco{O arco}
A série começou querendo detectar painéis solares por satélite e verificar
programas de crédito rural. Terminou medindo uma decisão operativa: quando, quanto
e por que o Operador Nacional do Sistema corta geração fotovoltaica --- e por que a
geração verificada deixou de ser proxy do recurso solar.

Três coisas mudaram o rumo. A \textbf{Fig.\,1} mostrou que o corte é sistêmico e de
razão energética, contaminando qualquer alvo de modelagem. A \textbf{Fig.\,2}
mediu que 91\% da micro e minigeração distribuída é sub-pixel no Sentinel-2,
fechando a rota original de auditoria por imagem. E a \textbf{Fig.\,8} entrou no
despacho, onde estava a explicação.

\bloco{Cinco figuras voltadas para dentro}
Cinco das doze não abrem tema novo: examinam o que a série já havia afirmado. Duas
\textbf{corrigem} --- derrubam uma conclusão anterior --- e três \textbf{delimitam}:
a Fig.\,10 responde ao nulo da Fig.\,3, a Fig.\,11 restringe o escopo da Fig.\,10 e a
Fig.\,12 mede a robustez das duas. As correções estão registradas em vez de apagadas.

A \textbf{Fig.\,6} derruba a conclusão secundária da Fig.\,5: o viés que atribuí à
referência do ONS era artefato do meu contrafactual --- 2023 já era um ano
restrito, apenas anterior ao início da apuração pública.

A \textbf{Fig.\,9} retira a atribuição causal da Fig.\,8: decomposta por
subsistema, a coincidência entre corte e térmica inflexível não sustenta leitura
local em nenhum dos dois subsistemas com corte.

\end{minipage}\hfill
\begin{minipage}[t]{0.47\linewidth}\vspace{0pt}
\bloco{Conteúdo}
\footnotesize
\begin{tabular}{@{}p{11mm}p{58mm}p{28mm}@{}}
\textbf{Método} & Janela, amostra, teste e multiplicidade & \textit{referência} \\
\textbf{Fig.\,1} & Corte fotovoltaico no SIN & mecanismo \\
\textbf{Fig.\,2} & MMGD: detectabilidade por satélite & viabilidade \\
\textbf{Fig.\,3} & Estudo de caso pareado, PIE $\times$ APE & caso \\
\textbf{Fig.\,4} & Cronologia do entorno & hipótese temporal \\
\textbf{Fig.\,5} & Controle climático & restrição ou recurso \\
\textbf{Fig.\,6} & A referência do ONS é enviesada? & \textit{corrige a 5} \\
\textbf{Fig.\,7} & Detecção de construção por Sentinel-2 & executa $H_2$ \\
\textbf{Fig.\,8} & Despacho fora da ordem de mérito & mecanismo \\
\textbf{Fig.\,9} & Decomposição por subsistema & \textit{corrige a 8} \\
\textbf{Fig.\,10} & Topologia da rede de transmissão & \textit{responde à 3} \\
\textbf{Fig.\,11} & Replicação em eólica & \textit{delimita a 10} \\
\textbf{Fig.\,12} & Robustez dos achados de rede & \textit{qualifica 10 e 11} \\
\textbf{Aux} & Imagem comercial: o que dá para comprar & decisão \\
\textbf{Aplic.} & E se eu quisesse construir ao lado? & exercício \\
\end{tabular}
\normalsize

\bloco{Fontes}
Exclusivamente dados abertos, sem autenticação: ONS (\textit{constrained-off},
geração térmica por motivo de despacho, fator de capacidade, curva de carga,
intercâmbio, energia armazenada, CMO), ANEEL (SIGA e geração distribuída),
Sentinel-2 L2A pelo catálogo STAC público, NASA POWER e malhas do IBGE.

\end{minipage}

\clearpage

%================================================================ MÉTODO
\cab{MÉTODO}{Janela, amostra e multiplicidade}{O recorte de cada figura, a unidade que se repete dentro dele, o teste que ela declara, o contrato de exportação e quantos testes a série reporta ao todo.}

\begin{minipage}[t]{0.635\linewidth}\vspace{0pt}
\bloco{Janela, amostra, unidade de repetição e teste}
\par\vspace{3pt}\small
\begin{tabular}{@{}lp{41mm}p{26mm}p{31mm}p{49mm}@{}}
\toprule
 & janela & $n$ & unidade de repetição & teste declarado \\
\midrule
\textbf{1} & abr/2024--jul/2026 & 85 usinas & usina · 30 min & nenhum: descritiva \\
\textbf{2} & conexões até abr/2026 & 4,6 M instalações & empreendimento & nenhum: descritiva \\
\textbf{3} & out/2024--ago/2026 (par) & 2 ativos; 60 conj. & conjunto & Spearman; 5 covariáveis, sem correção \\
\textbf{4} & 2023--jul/2026; entorno 2021 & 43 e 24 meses & conjunto · mês & nenhum teste formal \\
\textbf{5} & 2023--jul/2026 & meses com clima & conjunto · mês & nenhum: decomposição aritmética \\
\textbf{6} & abr/2024--jul/2026 & 74 calibradas & usina · dia-usina & Spearman; 2 covariáveis \\
\textbf{7} & 2018--2026 & 34 conj., 306 med. & conjunto · ano & Wilcoxon pareado; 2 comparações \\
\textbf{8} & abr/2024--ago/2026 & 21.192 horas & hora do SIN & Mann-Whitney e Spearman; quintis \\
\textbf{9} & idem, 9--15\,h & 21.192 horas & hora do subsistema & Spearman; sem correção \\
\textbf{10} & set/2025--ago/2026 & 59 conjuntos & conjunto & Spearman; 8 covariáveis, \textbf{Holm} \\
\textbf{11} & nov/2022--jul/2025 eól. & 120 conjuntos & conjunto & idem, \textbf{Holm} por fonte \\
 & set/2025--ago/2026 sol. & 59 conjuntos & conjunto &  \\
\textbf{12} & amostras das Figs.\,10 e 11 & 59 e 120 & conjunto & bootstrap e $z$ de Fisher \\
\bottomrule
\end{tabular}
\par\vspace{3pt}
\ressalva{As janelas de corte saem da mesma regra --- percentil 80 do início e 20 do fim da apuração --- e por isso diferem entre fontes. A Fig.\,1 vem de extração anterior às demais. Cada figura repete estas declarações na sua própria página, com as exclusões.}
\end{minipage}\hfill
\begin{minipage}[t]{0.335\linewidth}\vspace{0pt}\small
\bloco{Padrão de figura}
Todas as figuras seguem o mesmo contrato: 183\,mm de largura, corpo 7\,pt, piso de
5\,pt para qualquer glifo, texto editável em PDF e SVG, TIFF a 600\,dpi, dados-fonte
exportados em CSV e auditoria automática de colisão obrigatória. Cada figura declara
$n$, unidade de repetição, teste, exclusões e limites de interpretação.

\bloco{Multiplicidade, contada de uma vez}
A correção é aplicada \textbf{dentro} de figura --- Holm sobre as oito covariáveis
pré-declaradas nas Figs.\,10 e 11 --- e declarada ausente na Fig.\,3, onde as cinco
testadas deram todas nulas. \textbf{Não há correção entre figuras.} Somados os testes que
vão a arquivo a série reporta \textbf{49} --- 5 na Fig.\,3, 2 na 7, 5 na 8, 14 na 9, 8 na
10, 8 na 11 e 7 na 12 ---, dos quais \textbf{16 passam por Holm e 33 não}. Não é garimpo:
as oito covariáveis foram pré-declaradas e testadas na mesma ordem nas duas fontes, e a
Fig.\,12 audita sobreviventes. Mas quem quiser um limiar de série aplica por conta
própria --- este documento não o aplica.
\end{minipage}
\clearpage

%================================================================ FIG 1
\analise{ANÁLISE 1 \quad · \quad ONS, restrição por constrained-off}
{O corte fotovoltaico no SIN é sistêmico}
{Ele incide sobre toda a janela solar, cresce em ritmo semelhante no Nordeste e no Sudeste/CO, e é majoritariamente de razão energética --- excesso de oferta, não falha local.}
{../figures/fig1_curtailment_sin.pdf}
{
\bloco{O que a figura mostra}
\textbf{a}, perfil diurno da potência das usinas com apuração: azul é geração
verificada, vermelho é energia cortada. \textbf{b}, taxa mensal por subsistema.
\textbf{c}, energia cortada por razão e origem da restrição. \textbf{d},
distribuição espacial, cor pela taxa e tamanho pela capacidade.

\bloco{Números}
Corte de \textbf{21,6\%} no parque exposto. A razão energética domina; a origem é
predominantemente sistêmica.

\bloco{Uma armadilha de campo}
O campo \texttt{val\_geracaolimitada} \textit{não} é a energia cortada --- é o teto
ao qual a geração foi limitada. Sua correlação com a geração verificada é 0,86 e com
a perda real é \textbf{0,02}. Usá-lo superestimaria a perda grosseiramente. O corte
é
\[ \max\bigl(\text{referência} - \text{verificada},\; 0\bigr) \]
apenas nos intervalos com razão apurada.

\bloco{Por que isso reorienta tudo}
Se a geração de uma usina é limitada administrativamente, qualquer modelo que tente
prever geração a partir de irradiância --- ou de imagem de satélite --- está
ajustando um alvo que embute uma decisão operativa, não apenas física.

\lit{Os 21,6\% medidos aqui são de solar centralizada com apuração, na janela
abr/2024--jul/2026. O ONS reporta 20,7\% para eólica \textit{e} solar em 2025 --- ordens
de grandeza compatíveis, denominadores diferentes. No cenário internacional é um valor
alto: Chile perdeu 17\% da renovável variável disponível em 2024, Espanha bateu recorde
com $\approx$11\% em jul/2025, e na Califórnia a curtailment \textit{média} da solar é
4,3\%. Novan e Wang (2024) mostram que a taxa \textbf{marginal} é cerca de duas vezes a
média --- o que sugere que a exposição de uma usina nova no Brasil é ainda maior que
21,6\%.}

\bloco{Ressalvas}
\ressalva{O denominador é o das usinas com apuração, não do parque solar nacional.}
\ressalva{Não cobre MMGD, que tem regime de restrição distinto.}
\ressalva{A referência é estimativa do operador: em 21,8\% dos intervalos restritos ela fica abaixo da geração verificada, e o truncamento em zero torna a estimativa conservadora, não exata.}
\ressalva{Razão e origem são classificações operativas do ONS, não inferência causal independente.}
}{\legI}

%================================================================ FIG 2
\analisec{ANÁLISE 2 \quad · \quad ANEEL, geração distribuída}
{Nove em cada dez telhados são sub-pixel}
{A rota original --- verificar instalações registradas por imagem de satélite --- é inexequível com sensor gratuito, e o evento de instalação não é exógeno.}
{../figures/fig2_mmgd_detectabilidade.pdf}
{
\bloco{O teste}
A ANEEL publica a \textbf{área do arranjo} de cada instalação. Sobre 4,6 milhões de
registros, o número de pixels que cada sensor entrega é
\[ \mathrm{px}(s) = A_{\text{arranjo}} / r_s^2 \]
Isso transforma uma ressalva qualitativa em medida.

\bloco{Números}
Mediana de pixels por instalação e fração sub-pixel: Sentinel-2 (10\,m)
\textbf{0,29\,px} e \textbf{91,2\%}; CBERS-4A WPM (2\,m) 7,25\,px e 0,5\%; VHR
comercial (0,5\,m) 116\,px e $\approx$0.

A classe rural --- alvo dos programas de crédito --- tem área mediana de 41\,m$^2$ e
apenas \textbf{5,2\%} de segmentabilidade no Sentinel-2.

\bloco{O gargalo não é o sensor}
O registro \textbf{não tem localização utilizável}: os campos de coordenada existem
no esquema e estão preenchidos em \textbf{0,06\%}. Há \texttt{CodCEP} --- completo
para os 7,7\% de pessoa jurídica, truncado em 5 dígitos para os 92,3\% de pessoa
física --- mas ele \textbf{não refina o alvo rural}, com mediana de 1 prefixo por
município. Sem saber qual telhado olhar, verificar exige varredura exaustiva. Só o
segmento de pessoa jurídica é endereçável.

\bloco{E o antes/depois também cai}
A janela de transição da Lei 14.300/2022 concentrou 98 mil conexões num único mês ---
quatro vezes a média do ano anterior. A data de conexão não é exógena: um estudo de evento ingênuo mediria a
corrida regulatória, não a política.

\lit{A literatura que detecta telhados com sucesso opera em 0,1--0,3\,m: DeepSolar
(Yu et al., 2018) usa imagem aérea submétrica; o conjunto multirresolução de Jiang et al.
(2021) vai de 0,1 a 0,8\,m. O inventário global de Kruitwagen et al.\ (2021), em
\textit{Nature}, é o trabalho de referência com Sentinel-2 --- e cobre apenas instalações
acima de 10\,kW. Nenhum trabalho segmenta telhado residencial a 10\,m. A medida da
Fig.\,2 quantifica por quê.}

\bloco{Ressalvas}
\ressalva{Contar pixels não é detectar: um objeto sub-pixel ainda altera a assinatura espectral. O que a contagem descarta é a segmentação.}
\ressalva{O piso de 4 pixels é convenção declarada, não constante física.}
\ressalva{É censo administrativo, não amostra: não há erro amostral a declarar.}
}{\legII}

%================================================================ FIG 3
\analisec{ANÁLISE 3 \quad · \quad ONS + ANEEL/SIGA, estudo de caso pareado}
{Dois ativos, três vezes de diferença, nenhuma explicação que generalize}
{Sol do Cerrado perde 3,0$\times$ mais energia que Três Marias 3 --- mas nenhuma covariável local reproduz esse padrão no parque inteiro.}
{../figures/fig3_estudo_caso.pdf}
{
\bloco{O par}
Sol do Cerrado (681,3\,MW, produtor independente, 17 SPEs de um grupo, barramento
de Jaíba 230\,kV \textit{dividido}) contra Conj.\ Três Marias 3 (70,0\,MW,
autoprodução, sete empresas privadas distintas, ponto \textit{exclusivo}). Ambos em
Minas, mesmo subsistema, janela comum recortada programaticamente.

\bloco{Uma correção de atribuição}
O campo \texttt{nom\_agenteproprietario} do ONS identifica o \textbf{agente
representante}, não o proprietário. Foi o que sugeriu, incorretamente, que Três
Marias 3 seria ativo estatal da CEMIG. O SIGA mostra sete proprietários privados sob
regime de autoprodução --- não existe usina solar estatal relevante no parque
despachado, então o eixo público $\times$ privado não é testável.

\bloco{Números}
Taxa de corte 26,8\% contra 8,8\%. A energia cortada vale cerca de
\textbf{R\$\,64 milhões} contra pouco mais de R\$\,1 milhão, a CMO.

\bloco{O achado de despacho, ainda sem nome}
Nos dois ativos o CMO ponderado \textbf{durante o corte} é bem menor que durante a
geração. O sistema corta quando a energia vale pouco --- coerente com restrição de
razão energética. Este é o fio que a Fig.\,8 puxa.

\bloco{O resultado negativo}
Duas explicações locais eram plausíveis: barramento compartilhado e densidade
regional de renovável. Nenhuma sobrevive ao teste na frota --- nem carga do ponto de
conexão, nem número de conjuntos, nem densidade regional, nem tensão, nem capacidade
própria. \textbf{Explicações locais convincentes não devem ser promovidas a mecanismo
geral sem esse teste.}

\bloco{Uma correção sobre o que é público}
A versão original desta análise concluía que a topologia de transmissão não está em
dado aberto. \textbf{Está.} O SIGET publica o grafo linha a linha --- 1.166 linhas,
619 subestações com código ONS, tensão e comprimento --- e o \texttt{fator-capaci\-dade-2}
traz a coordenada do ponto de conexão de cada usina. O join cobre 168 dos 229
conjuntos despachados. E o que eu disse faltar --- o \textbf{limite de transferência}
--- também está publicado: \texttt{linha-transmissao} do ONS traz capacidade operativa
em MVA com variantes sazonais, mais resistência e reatância. Falta apenas
\textbf{qual restrição está ativa em cada intervalo}. As cinco covariáveis testadas
são locais ao ponto e não usam a estrutura da rede; o nulo não se estende a ela --- e
a \textbf{Fig.\,10} a testa, achando duas associações que sobrevivem à correção de
multiplicidade.

\lit{Newbery (2024): quem é cortado depende do regime de acesso, não só da física da rede.}

\bloco{Ressalvas}
\ressalva{Dois ativos não são amostra: o par foi escolhido por contraste.}
\ressalva{Cinco covariáveis testadas sem correção de multiplicidade; como nenhuma resultou significativa, a correção só tornaria a conclusão mais conservadora.}
\ressalva{CMO é proxy de valor sistêmico, não preço de liquidação nem prejuízo líquido.}
}{\legIII}

%================================================================ FIG 4
\analise{ANÁLISE 4 \quad · \quad ONS, fator de capacidade}
{O corte acompanha a energização dos vizinhos?}
{A apuração de corte começa em abr/2024 e a onda de construção em Jaíba é de 2021 a 2024. O fator de capacidade contorna a janela.}
{../figures/fig4_cronologia.pdf}
{
\bloco{O problema de janela}
Olhar apenas a série de corte deixaria de fora justamente o período em que a
hipótese deveria se manifestar. O fator de capacidade existe desde a entrada em
operação e cai quando a geração é limitada:
\[ FC_{\text{verificado}}(t)=\frac{E_{\text{gerada}}(t)}{P_{\text{inst}}\cdot\Delta t} \]

\bloco{A decomposição}
A partir de abr/2024 a apuração permite recompor o potencial:
\[ FC_{\text{potencial}} = FC_{\text{verificado}} + FC_{\text{cortado}} \]
Se a queda for de restrição de rede, o potencial recomposto permanece estável
enquanto o verificado cai. Se for recurso solar ou degradação, os dois caem juntos.

\bloco{O que isso deixa em aberto}
A decomposição separa restrição de degradação, mas não separa restrição de
\textbf{variação climática} --- se 2025 tiver sido menos ensolarado que 2023, a
queda não diz nada sobre a rede. Essa lacuna é o que a Fig.\,5 fecha, medindo a
irradiância diretamente sobre cada usina.

\lit{O desenho é um estudo de evento com adoção escalonada --- cada vizinho energiza numa
data diferente --- e é exatamente o caso em que Callaway e Sant'Anna (2021) mostram que o
estimador de diferenças em diferenças com efeitos fixos duplos fica enviesado, por ponderar
comparações entre unidades já tratadas. A leitura desta figura é descritiva e pareada
contra um único controle, o que evita o problema mas também não estima efeito de
tratamento. Estimá-lo exigiria o estimador por grupo e período daqueles autores --- e uma
definição de tratamento exógena, que a Fig.\,2 mostrou não existir do lado da MMGD.}

\bloco{Ressalvas}
\ressalva{O fator de capacidade é observável de longo prazo, mas responde a várias causas ao mesmo tempo; a decomposição só vale onde a apuração existe.}
\ressalva{Comparação pareada sem aleatorização, como na Fig.\,3.}
}{\legIV}

%================================================================ FIG 5
\analise{ANÁLISE 5 \quad · \quad ONS + NASA POWER}
{A queda é restrição, não recurso}
{Normalizada pela irradiância medida sobre a usina, a perda de Sol do Cerrado permanece: nove décimos da queda não são atribuíveis ao sol.}
{../figures/fig5_controle_climatico.pdf}
{
\bloco{O método}
Com a irradiância global horizontal diária $H$ sobre a usina, a razão de desempenho
\[ PR = \frac{FC \cdot 24\,\mathrm{h}}{H} \]
divide a energia gerada por kW instalado pelo recurso disponível. Se a queda fosse
de menos sol, $PR$ ficaria estável.

\bloco{Números}
A irradiância sobre Jaíba caiu \textbf{2,8\%} entre 2023 e 2025. O fator de
capacidade caiu \textbf{29,1\%}. Escalando o FC de 2023 pela razão de irradiância, o
esperado para 2025 seria 0,206; o observado foi 0,150. Cerca de um décimo da queda é
do recurso; nove décimos não são.

A razão de desempenho cai de 0,854 para 0,619 em Sol do Cerrado. No controle, Três
Marias 3 fica em torno de 0,99 sem tendência.

\bloco{A relação se deslocou, não percorreu}
Os meses de 2025 não estão na mesma curva de 2023 em posição de menor irradiância
--- estão numa curva inteiramente mais baixa. É a assinatura de uma perda que
independe do recurso.

\bloco{E uma conclusão secundária que estava errada}
O déficit não explicado pelo sol é de 0,056 em fator de capacidade; o corte apurado
pelo ONS equivale a 0,147 --- 2,6 vezes maior. Li isso como indício de que a
referência do ONS superestima o contrafactual. \textbf{A leitura tinha um
pressuposto não testado: que 2023 fosse um ano livre de restrição.} A Fig.\,6
mostra que não era.

\lit{A irradiância aqui é variável derivada de sensoriamento, e entra como denominador da
razão de desempenho. Proctor, Carleton e Sum (2023) mostram que usar uma variável assim
numa inferência a jusante enviesa o ponto e o erro padrão mesmo com erro de medida
pequeno, e recomendam tratá-la como distribuição --- imputação múltipla --- e não como
valor pontual. Alix-García e Millimet (2023) acrescentam que o erro de produto de satélite
raramente é clássico. Nenhuma das duas correções foi aplicada: a diferença medida é grande
o bastante para sobreviver a erro moderado, mas o intervalo em torno dela é mais estreito
do que deveria ser.}

\bloco{Ressalvas}
\ressalva{Reanálise, não piranômetro: NASA POWER é MERRA-2 combinado com satélite, em célula de 0,5$^\circ\times$0,625$^\circ$.}
\ressalva{GHI no lugar de POA: a irradiância é horizontal, os módulos não. O PR aqui é indicador relativo dentro de cada usina, não PR de norma IEC.}
\ressalva{Temperatura não modelada; contrafactual linear.}
}{\legV}

%================================================================ FIG 6
\analise{ANÁLISE 6 \quad · \quad correção da Fig.\,5}
{A referência do ONS passa no teste --- o viés era meu}
{Em 74 usinas o viés mediano é 1,02. O déficit aparente da Fig.\,5 vinha do contrafactual, não do operador: 2023 já era um ano restrito.}
{../figures/fig6_vies_referencia.pdf}
{
\bloco{O teste}
Para cada usina, calibra-se uma razão de desempenho \textbf{nos dias sem restrição
da própria usina} ($f<1\%$) e compara-se o potencial implicado pelo ONS
(verificado $+$ cortado) com o potencial previsto pela irradiância local.

\bloco{Resultado}
Viés mediano por usina \textbf{1,02}, metade entre 0,98 e 1,06. Nenhuma das 74
usinas calibradas passa de 1,5, e 26 ficam abaixo de 1 --- para um terço delas a
referência do ONS é \textit{conservadora}. No agregado energético dos dias
restritos, o corte apurado é 1,03 vezes o implicado pela irradiância.

E não piora onde o corte é mais intenso, que seria a falha mais preocupante: a
associação entre viés e intensidade de corte é nula.

\bloco{A correção}
O desempenho limpo de Sol do Cerrado, calibrado em 279 dias sem restrição, é
\textbf{0,931}. Em 2023 a usina operou a 0,854 --- \textbf{8,3\% abaixo do seu
próprio desempenho limpo}. Aquele ano já era restrito; simplesmente não havia
apuração pública. Foi por isso que a Fig.\,5, ao usar 2023 como base, mediu um
déficit pequeno e concluiu, erradamente, que o ONS inflava o corte.

\bloco{O que muda e o que não muda}
A conclusão central da Fig.\,5 --- perda de rede, não de clima --- permanece e sai
reforçada: contra o desempenho limpo, o corte implícito é \textbf{33\% em 2025}, não
26\%. Cai a afirmação secundária sobre viés do operador. E os 21,6\% da Fig.\,1 não
precisam de correção: se a referência é calibrada, o agregado é confiável.

\lit{A geração de referência do ONS é derivada de curvas de desempenho da usina
(irradiância ou vento contra potência), como descrevem Vieira et al.\ (2026). Validá-la
por fora exige uma fonte independente: no lado eólico, Silva et al.\ (2026) fazem isso
com anemometria de SCADA. Este notebook faz o equivalente no lado solar, com irradiância
de satélite --- e é, até onde encontrei, a primeira verificação pública da calibração
dessa referência.}

\bloco{Ressalvas}
\ressalva{Dias ``limpos'' podem não ser limpos; como o resultado é ausência de viés, esse erro só o tornaria mais forte.}
\ressalva{Seis usinas ficaram de fora por não atingirem 30 dias limpos --- justamente as mais restritas, onde a referência é mais usada.}
\ressalva{O $PR$ limpo é constante por usina e herda qualquer viés sazonal dos dias de calibração.}
}{\legVI}

%================================================================ FIG 7
\analisec{ANÁLISE 7 \quad · \quad Sentinel-2 L2A (STAC público) + ANEEL/SIGA}
{O satélite vê a obra, não o painel}
{A construção de uma usina fotovoltaica tem assinatura robusta no Sentinel-2 e data a entrada em operação; a operação em si não se separa do entorno a 10\,m.}
{../figures/fig7_deteccao_construcao.pdf}
{
\bloco{O desenho}
Para cada conjunto, mede-se anualmente a razão entre o brilho no vermelho dentro do
arranjo e num anel de controle:
\[ \rho(t)=\frac{\overline{B_{\text{red}}}\big|_{d\le 400\,\mathrm{m}}}
{\overline{B_{\text{red}}}\big|_{1{,}5\le d\le 3\,\mathrm{km}}} \]
O anel controla iluminação, estação e umidade do solo. As séries são alinhadas pela
\textbf{data de entrada em operação} do SIGA, não pelo calendário. 34 conjuntos,
9 anos, 306 medições.

\bloco{A obra tem assinatura forte}
A razão sobe de 1,01 na base para 1,30 na obra --- \textbf{+34\% na mediana, 24 de 30
usinas}, Wilcoxon $p=1{,}2\times10^{-5}$. Construir exige remover vegetação de
centenas de hectares, e solo exposto é muito mais claro no vermelho.

\bloco{A operação não tem}
Da base à operação a mediana vai de 0,95 para 1,02 --- \textbf{$+6{,}6\%$}, com 11 de
18 usinas em alta contra 7 em queda ($p=0{,}09$). Não é escurecimento: é ausência de
sinal em qualquer direção. Módulos são escuros, mas a 10\,m o pixel mistura painel,
corredor, via e solo.

\bloco{Uma correção de método que domina o resultado}
A primeira tentativa usou a coordenada do \textbf{conjunto} publicada pelo ONS, que
é a da subestação coletora. O resultado foi nulo. Em \textbf{31 de 37} conjuntos a
subestação fica a mais de 750\,m do centroide do arranjo --- mediana de 1,9\,km. A
máscara media o chão errado. Refeito com as coordenadas de usina do SIGA, o sinal
aparece.

\bloco{O que responde da proposta original}
A hipótese ``registro $\leftrightarrow$ evidência no satélite'' é executável para
geração centralizada, com uma inversão: o satélite confirma a usina pela
\textbf{cicatriz da obra}, não pelo painel. É evento transitório --- para auditoria,
é preciso monitorar continuamente, não fotografar o estado atual.

\lit{A razão dentro/fora do arranjo é, em essência, uma versão simplificada do
\textit{Temporal Cluster Matching} (Robinson et al., Microsoft Research), que data
construções comparando a relação espectral dentro e fora da pegada --- e foi aplicado
justamente a usinas solares em Karnataka. O Global Renewables Watch estima datas de
construção em escala global por imagem de alta resolução. O que a Fig.\,7 acrescenta é a
validação contra o \textbf{registro de outorga}, e o achado negativo de que a fase
operacional não é separável a 10\,m.}

\bloco{Ressalvas}
\ressalva{A qualidade da geolocalização de referência domina o resultado; qualquer replicação depende dessa camada.}
\ressalva{Uma cena por ano fixa a datação em resolução anual: acerto exato em 21\% e dentro de um ano em 62\% dos casos.}
\ressalva{Círculos de 400\,m não são o polígono do arranjo; delinear por segmentação melhoraria o contraste.}
\ressalva{O pulso indica remoção de vegetação, que também ocorre em lavoura, mineração e loteamento: o que valida a atribuição é o cruzamento com o registro. E nada disso se estende a telhados.}
}{\legVII}

%================================================================ FIG 8
\analise{ANÁLISE 8 \quad · \quad ONS, geração térmica por motivo de despacho}
{O corte por excesso de oferta coincide com térmica inflexível}
{Dois terços do corte são classificados como excesso de oferta e ocorrem com 4,6\,GWmed de térmica fora da ordem de mérito no mesmo instante, 73\% dela inflexibilidade declarada.}
{../figures/fig8_despacho_merito.pdf}
{
\bloco{A definição que fecha}
A decomposição do ONS é \textbf{aninhada, não aditiva}: somar todos os campos dá
1,58 vez a geração total. A única definição consistente é
\[ \text{fora do mérito} = \texttt{verifgeracao} - \texttt{verifordemmerito} \]
verificada contra a soma dos motivos não-econômicos disjuntos --- razão
\textbf{1,0030}, mediana da diferença 0,17\,MWmed.

\bloco{A leitura tentadora, e por que erra}
``A térmica não abre espaço para a solar'' não sobrevive ao controle de carga: a
carga sobe quase o quanto a solar sobe, e a \textbf{carga líquida é plana}
(72,2 $\rightarrow$ 72,1\,GWmed).

\bloco{O que muda é a classificação}
Da madrugada ao pico solar a parcela por mérito cai \textbf{1.640\,MWmed} e a
parcela fora do mérito sobe \textbf{1.184\,MWmed}. A mesma térmica fica ligada e
deixa de ser econômica: quando o CMO desaba, mérito vira inflexibilidade.

\bloco{Números}
21.889\,GWh de corte na janela, 68,5\% por razão energética. Nas horas com esse
corte há 4.564\,MWmed de térmica fora do mérito contra 3.146 nas demais (+45\%);
ponderado pela energia cortada, 5.266. \textbf{73\% é inflexibilidade pura.}

\bloco{E é condicional}
A correlação entre corte e térmica inflexível é $\rho=0{,}48$ no quintil de carga
mais baixa, 0,38 no segundo, e cai para 0,09--0,20 nos três superiores. Com $n\approx1.240$ por estrato
\textbf{todos os quintis dão $p<0{,}05$} --- a leitura se apoia na magnitude, não na
significância, e a queda não é monotônica.

\lit{O mecanismo é canônico na literatura chinesa de curtailment: unidades térmicas
com mínimo técnico e status \textit{must-run} ocupam a faixa que a renovável variável
precisaria, e a resposta de política foi o programa de \textit{flexibility retrofit} ---
16,35\,GW em demonstração em 2017, somando 4,8\,GW de capacidade de regulação de ponta.
No Brasil, Vieira et al.\ (2026) modelaram o corte com regressão logística horária e
acharam a MMGD como preditor mais forte do corte por demanda ($OR\approx9{,}8$), com a
carga do sistema atuando de forma protetiva --- exatamente o gradiente que o painel
\textbf{c} reproduz.}

\bloco{Ressalvas}
\ressalva{Coincidência não é causalidade: separar exigiria reotimizar o despacho sem a inflexibilidade declarada.}
\ressalva{Inflexibilidade é declaração do agente, não mínimo técnico verificável em dado aberto.}
\ressalva{Agregado nacional --- corrigido na Fig.\,9.}
\ressalva{Controle negativo: energia armazenada não explica ($\rho=0{,}06$, $p=0{,}06$).}
}{\legVIII}

%================================================================ FIG 9
\analise{ANÁLISE 9 \quad · \quad correção da Fig.\,8}
{O piso térmico está onde o corte não está}
{Decomposta por subsistema, a coincidência nacional não sustenta leitura causal local em nenhum dos dois subsistemas com corte.}
{../figures/fig9_decomposicao_subsistema.pdf}
{
\bloco{O descompasso}
Nordeste e Sudeste cortam volumes parecidos nas horas solares. A térmica inflexível
não acompanha.
\par\vspace{1pt}\scriptsize
\begin{tabular}{@{}lrr@{}}
\toprule
MWmed, horas de 9\,h a 15\,h & Nordeste & Sudeste/CO \\
\midrule
corte por razão energética & 1.079 & 950 \\
inflexibilidade pura local & \textbf{36} & 1.739 \\
corte / piso local & \textbf{30,2$\times$} & 0,5$\times$ \\
$\rho$ com térmica do NE & +0,08 & +0,10 \\
$\rho$ com térmica do SE & \textbf{+0,42} & \textbf{+0,41} \\
\bottomrule
\end{tabular}
\par

\bloco{No Nordeste o mecanismo está eliminado}
Onde há o maior volume de corte do país há 36\,MWmed de inflexibilidade contra
1.079 de corte. A associação local é $\rho=0{,}08$ e não passa de 0,16 em nenhum
quintil de carga.

\bloco{E o teste atinge também o Sudeste}
Eu esperava que o mecanismo sobrevivesse ali, e o painel \textbf{b} parecia dar isso:
$\rho=0{,}40$ nos quintis de carga baixa, o padrão exato da Fig.\,8. O painel
\textbf{c} desfaz a inferência. A térmica do Sudeste se correlaciona com o corte do
Sudeste em 0,41 e com o do \textbf{Nordeste} em 0,42, a mais de dois mil
quilômetros. \textbf{O que determina o $\rho$ é qual térmica se olha, não qual
corte} --- e uma causa local produziria associação mais forte com o corte local.

\bloco{A hipótese de corredor também cai}
Nas horas de corte alto o fluxo NE$\rightarrow$SE é \textit{menor} --- 2.993 contra
3.848\,MWmed --- e fica acima do teto empírico em 2,9\% das horas contra 5,7\% nas
demais ($\rho=-0{,}25$). O Nordeste é cortado com o corredor \textbf{vazio}.

\bloco{Efeito sobre a Fig.\,8}
Nenhum número dela muda. O que se retira é a \textbf{atribuição causal}: a
coincidência é real e compatível com um fator comum --- carga nacional baixa --- e
este desenho não distingue as duas hipóteses.

\lit{Newbery (2024) trata precisamente de \textit{quem} é cortado sob restrição de
transmissão, e de como o regime de acesso distorce o sinal de investimento --- o
enquadramento econômico do descompasso medido aqui. E a invisibilidade da MMGD ao
despacho é problema documentado: geração atrás do medidor não é observada pelo operador
e eleva o erro de previsão de carga líquida em 30\% no nível domiciliar e até 250\% no
agregado.}

\bloco{Ressalvas}
\ressalva{Subsistema ainda é agregação grosseira: o Sudeste/CO vai de Minas ao Mato Grosso do Sul.}
\ressalva{O teto do corredor é o percentil 95 do próprio fluxo. Verificado: \texttt{programacao\_fluxo\_controlado} cobre só HVDC e interligações internacionais, e \texttt{linha-transmissao} dá capacidade por linha, não do corredor --- que resulta de estudo de estabilidade.}
\ressalva{``Ausência de demanda'' é inferência por eliminação, não testada positivamente.}
\ressalva{O painel c não prova ausência de mecanismo local no Sudeste --- mostra que este desenho não o identifica, porque inflexibilidade local e carga nacional são colineares.}
}{\legIX}

%================================================================ FIG 10
\analisea{ANÁLISE 10 \quad · \quad EPE (rede) + ONS \quad · \quad responde à Fig.\,3}
{A distância até a demanda explica parte do corte}
{Com o grafo de transmissão montado, duas de oito covariáveis pré-declaradas sobrevivem à correção de multiplicidade --- o primeiro achado positivo da série sobre a alocação do corte.}
{../figures/fig10_topologia_rede.pdf}
{
\bloco{O que mudou em relação à Fig.\,3}
Aquela figura testou cinco covariáveis \textbf{locais ao ponto de conexão} e não achou
nenhuma. Aqui elas são \textbf{estruturais} --- só existem se houver um grafo. Rede da
EPE: 920 subestações, 1.896 linhas resolvidas em aresta (87\%), ancoragem mediana de
0,10\,km.

\bloco{Resultado}
Duas sobrevivem a Holm, com magnitude quase igual e sinais opostos:
\textbf{distância de rede à carga} ($+$0,38) e \textbf{capacidade da usina}
($-$0,38), ambas com Holm $p = 0{,}023$. A soma de tensões no nó chega a
$p < 0{,}05$ bruto e cai sob correção; as outras cinco não chegam.

\bloco{A contraparte transversal da Fig.\,9}
Quanto mais longe do centro de carga, em quilômetros de linha percorrida, maior o
corte. Na Fig.\,9 o Nordeste era cortado com o corredor NE$\rightarrow$SE
\textit{vazio}; aqui a usina mais cortada é a mais distante de quem consome. Duas
medidas independentes, mesmo mecanismo: \textbf{a distância até a demanda, não a
saturação local do nó}.

\bloco{E não é apenas subsistema}
Era o risco depois da Fig.\,9. A distância mediana à carga é quase igual nos dois
--- 497\,km no Nordeste, 520 no Sudeste/CO --- e a associação se mantém
\textbf{dentro} do Nordeste ($\rho = +0{,}34$, $p = 0{,}027$, $n = 43$), mas não dentro
do Sudeste/CO ($+0{,}31$, $p = 0{,}25$, $n = 16$). A capacidade da usina se mantém nos
\textit{dois} --- $-0{,}41$ e $-0{,}44$ ---, e no Nordeste com $p = 0{,}006$: mais firme
que a distância. A estrutura local do nó segue sem explicar nada.

\bloco{Uma sobrevivente desconfortável}
A capacidade da usina tem sinal \textbf{negativo}: as maiores são cortadas menos. Não
era hipótese direcional. As explicações plausíveis --- qualidade de conexão, seleção
contratual --- são testáveis e \textbf{não foram testadas}. E ela é, na Fig.\,12, o
\textbf{mais robusto dos dois achados solares}: passa nos dois testes que se aplicam a
ela, enquanto a distância à carga falha em dois dos três. Promovi a errada a principal.

\bloco{Escopo, delimitado duas vezes}
Vale para fotovoltaica: replicado em eólica \textbf{não se sustenta}, e a Fig.\,11
mostra por quê. E a Fig.\,12 mede sua robustez: intervalo bootstrap de largura 0,49,
coeficiente que cai de 0,62 para 0,11 conforme a metade da amostra, e ausência no
grafo construído por barra. \textbf{É o mais frágil dos quatro achados de rede da
série} --- e os testes que mostram isso poderiam ter sido feitos aqui.

\lit{Maji et al.\ (e-Energy 2025) defendem que o corte é fenômeno de \textbf{nó}, e que
estimá-lo exige descer ao ponto de conexão. É a hipótese que esta figura testa --- e que,
no fotovoltaico, ela \textit{não} confirma: o que prevê é a distância até a carga,
grandeza de caminho, não de nó. A Fig.\,11 mostra que em eólica vale o contrário, o que
delimita a afirmação daqueles autores por composição da causa. Newbery (2024) dá o
enquadramento do que isso implica: se a posição na rede determina a exposição ao corte e o
regime de acesso não a precifica, o sinal de investimento aponta para o lugar errado.}

\bloco{Ressalvas}
\ressalva{Distância de rede não é elétrica: caminho ponderado por km de linha. Eu havia escrito que impedância não estaria em dado aberto --- está, em 100\% das linhas de \texttt{linha-transmissao} (ONS), com capacidade em MVA em 86,7\%. A escolha da camada geométrica da EPE foi minha, e custou a distância elétrica.}
\ressalva{Centros de carga são nove metrópoles, proxy geográfico declarado.}
\ressalva{n = 59 para oito hipóteses. Holm foi aplicado, mas não há poder para efeitos pequenos. É correlação, não causa.}
}{\legX}

%================================================================ FIG 11
\analisec{ANÁLISE 11 \quad · \quad ONS (eólica) + EPE \quad · \quad delimita a Fig.\,10}
{A replicação falha --- e a falha explica o mecanismo}
{O achado da Fig.\,10 não se sustenta em eólica. Mas a covariável que prevê o corte de cada fonte espelha a origem declarada do corte dessa fonte.}
{../figures/fig11_replicacao_eolica.pdf}
{
\bloco{O teste mais duro disponível}
O ONS apura \textit{constrained-off} eólico desde \textbf{out/2021} em \textbf{199
conjuntos} contra 87 do solar. Mesmas oito covariáveis da Fig.\,10, mesma ordem, mesmo
grafo, Holm reaplicado.

\bloco{Não replica}
A distância de rede à carga, que no solar dava $\rho = +0{,}38$ e sobrevivia a Holm,
dá $\rho = -0{,}07$ em 120 conjuntos eólicos. Não é falta de amostra --- a eólica tem
o dobro.

E as covariáveis que funcionam são as que falharam no solar: usinas no mesmo nó vai de
$+0{,}07$ para \textbf{$+0{,}47$}, grau do nó de $-0{,}17$ para \textbf{$+0{,}35$},
ambas sobrevivendo a Holm com folga.

\bloco{A explicação está medida}
O ONS classifica cada intervalo restrito por origem, independentemente de tudo o que
foi testado. O corte solar é \textbf{91\% sistêmico}; o eólico, 76\% --- com
\textbf{24\% de origem local}, 2,7 vezes a fração do solar. Confiabilidade local é
7\% do corte solar e \textbf{22\%} do eólico.

\textbf{A covariável que prevê o corte de cada fonte espelha a origem declarada do
corte dessa fonte.} Sistêmico prevê-se por grandeza sistêmica --- distância até a
demanda; local, por grandeza local --- quantos disputam o ponto.

\bloco{Não é geografia}
A amostra eólica é 91\% nordestina e a solar não é. Restringindo as duas ao Nordeste,
o solar mantém distância à carga em $+0{,}34$ e usinas no nó em $+0{,}05$; a eólica
inverte, com $-0{,}19$ e $+0{,}36$. A dissociação sobrevive inteira.

\bloco{O que faz com a Fig.\,10}
Não a derruba --- no fotovoltaico o achado segue de pé. Muda a generalidade. Eu havia
escrito que a hipótese era ``sobre posição na rede, não sobre tecnologia''.
\textbf{Estava errado}: é sobre posição na rede \textbf{para fonte de corte sistêmico}.

\bloco{Consequência prática}
Solar pede demanda mais perto --- armazenamento, deslocamento de carga. Eólica pede
reforço local de escoamento. A mesma política erraria em uma delas.

\bloco{Robustez, medida na Fig.\,12}
\textbf{Grau do nó} passa nos três testes e é o achado mais firme da série sobre
alocação de corte. \textbf{Usinas no nó} replica no grafo alternativo mas é instável
entre subamostras: a direção se mantém, a magnitude não.

\lit{É o eixo canônico da literatura: Bird, Lew e Milligan (2016) organizam a revisão
de onze países por excesso de oferta sistêmico contra restrição local. A EirGrid
formaliza no vocabulário --- \textit{curtailment} para razão sistêmica,
\textit{constraint} para local --- e reporta 6,6\% contra 4,7\% no
\textit{dispatch-down} eólico de 2025. Maji et al.\ (e-Energy 2025) defendem que o
corte é fenômeno de nó. O passo que não encontrei publicado é o desta figura: testar
se a \textbf{covariável que prevê o corte muda com a composição da causa}.}

\bloco{Ressalvas}
\ressalva{As janelas não coincidem --- eólica de nov/2022 a jul/2025, solar de set/2025 a ago/2026. Parte da diferença pode ser de período; o teste no Nordeste controla geografia, não tempo.}
\ressalva{A eólica se aglomera mais (mediana de 4 por nó contra 2 no solar), o que facilita detectar aglomeração --- mas a distância à carga tem faixa semelhante nas duas e mesmo assim não prevê nada em eólica.}
\ressalva{Regimes distintos nesta amostra: 8\% mediano em eólica e 26\% em solar; nas frotas, 15,9\% e 22,1\%. Origem é classificação do ONS, não inferência causal.}
}{\legXI}

%================================================================ FIG 12
\analisec{ANÁLISE 12 \quad · \quad Figs.\,10 e 11 + ONS \quad · \quad qualifica ambas}
{Quanto vale cada achado}
{A única análise da série que não produz achado novo: submete os quatro resultados de rede que sobreviveram a Holm a testes que eu não tinha aplicado, e ordena o que sobra.}
{../figures/fig12_robustez.pdf}
{
\bloco{Por que esta figura existe}
Ao levantar fontes descobri que o ONS publica a rede com \textbf{conectividade por
número de barra}, impedância em 100\% das linhas e capacidade em MVA em 87\% ---
fonte melhor que a camada geométrica usada nas Figs.\,10 e 11. Reconstruí o grafo por
essa via, e o resultado obrigou a esta análise.

\bloco{O placar}
\par\vspace{1pt}\scriptsize
\begin{tabular}{@{}lccc@{}}
\toprule
 & IC & estável & outro grafo \\
\midrule
eólica · grau do nó & sim & sim & sim \\
eólica · usinas no nó & sim & \textbf{não} & sim \\
solar · capacidade da usina & sim & sim & --- \\
solar · distância à carga & sim & \textbf{não} & \textbf{não} \\
\bottomrule
\end{tabular}
\par
\ressalva{--- não se aplica: a capacidade da usina não vem do grafo. Dois critérios, e não três.}

\bloco{Quatro, e não três}
Sobreviveram a Holm \textbf{quatro} achados. A \textbf{capacidade da usina} em solar
--- $\rho = -0{,}38$, Holm $p = 0{,}023$, maior que a distância à carga --- ficara de
fora.

\bloco{O achado principal da Fig.\,10 é o mais frágil}
Intervalo bootstrap de 0,12 a 0,61 --- \textbf{largura 0,49}, quase o dobro dos
outros. Partindo a amostra, cai de \textbf{0,62} para \textbf{0,11} conforme a metade
($p = 0{,}025$), e a metade que o carrega contém \textbf{todas as 16 usinas do
Sudeste} --- o que reabre o confundimento regional que o painel \textbf{d} da Fig.\,10
afirmava ter afastado.

\bloco{E dentro do solar a ordem se inverte}
A capacidade da usina passa nos dois critérios que se aplicam: IC de $-0{,}60$ a
$-0{,}13$ e partição não rejeitada ($-0{,}18$ contra $-0{,}48$, $p = 0{,}21$). O que a
Fig.\,10 promoveu a principal é o que menos resiste. Mas dois testes não são três, e a
diferença entre metades ainda some dois terços: \textbf{passar aqui é não ser
reprovada, não é ser confirmada}.

\bloco{Amostra, não método}
No grafo do ONS a distância dá $-0{,}14$ em km e $-0{,}06$ em reatância. Mas rodando o
grafo \textit{antigo} nas mesmas 28 usinas, ele também não acha nada: a diferença é de
amostra, o que não salva o achado --- apenas identifica onde ele mora.

A distância elétrica foi calculada, e a Fig.\,10 errava ao dizer que não estaria
disponível. Sua correlação com a distância em km é de apenas \textbf{0,53}, então a
troca era teste real. Não mudou o resultado --- mas isso se sabe agora.

\bloco{A lição de método}
Os três testes \textbf{poderiam ter sido feitos na Fig.\,10}, com os dados que já
estavam na mesa. Não dependiam de nenhuma fonte nova. Foram feitos porque a troca de
grafo os forçou.

\lit{Proctor, Carleton e Sum (2023): variável derivada de sensoriamento --- como este
grafo --- propaga erro para o intervalo, e nenhuma imputação foi aplicada aqui.}

\bloco{Ressalvas}
\ressalva{A partição não é aleatória, de propósito: testa dependência a característica estrutural, não variabilidade amostral genérica.}
\ressalva{O grafo do ONS é exato onde cobre, mas exclui conexões abaixo de 230 kV --- para a amostra solar isso corta mais da metade e elimina um subsistema inteiro.}
\ressalva{Três testes não esgotam robustez: faltam validação temporal, análise de influência ponto a ponto e sensibilidade à definição de centro de carga. Os limiares do placar são convenção declarada, e o teste do IC foi corrigido de ``é positivo'' para ``exclui zero''.}
}{\legXII}

%================================================================ AUX
\analise{PÁGINA AUXILIAR \quad · \quad apoio de decisão, fora da série}
{O que dá para comprar de imagem comercial}
{Os 21 sensores de um fornecedor típico posicionados contra a curva de detectabilidade medida na Fig.\,2, e o preço do campo que a ANEEL não publica.}
{../figures/aux_geopera_capacidade.pdf}
{
\bloco{A leitura que o catálogo não mostra}
Toda a transição útil está \textbf{entre 1\,m e 5\,m}. Abaixo de 1\,m, 99,5\% da
MMGD é segmentável e os 18 sensores submétricos estão todos no platô --- resolução
extra não compra nada para este problema. Acima despenca: 86,6\% a 2\,m, 7,5\% a
5,3\,m, 1,6\% no Sentinel-2. Os dois sensores hiperespectrais, que são a parte mais
distintiva do catálogo, caem do lado errado do joelho.

\bloco{O tier barato já é gratuito}
A faixa de 2\,m custa AUD 2--3/km$^2$ e entrega a mesma resolução do CBERS-4A WPM,
que é público. De 0,80\,m a 2\,m só há arquivo --- não se pode tasquear.

\bloco{Cinco ordens de grandeza}
O painel \textbf{c} precifica a ressalva da Fig.\,2. Com coordenada por instalação,
um chip de 0,01\,km$^2$ a 30\,cm custa AUD 0,18--0,28. Sem ela, é preciso varrer o
município: 1.000\,km$^2$ custam AUD 18 mil--28 mil. \textbf{O gargalo do desenho não
é sensor nem orçamento --- é governança de dado.}

\bloco{O que isso destrava, e o que não}
Destrava uma versão reduzida: varrer 100--400\,km$^2$ da área ocupada de dois ou três
municípios rurais custa AUD 1.800--11.200 e permite auditar a \textbf{contagem
agregada} contra o registro municipal. Não destrava a verificação por instalação
\textit{na classe rural}.

Há, porém, uma exceção que o cálculo do painel \textbf{c} não captura: as 361.622
instalações de \textbf{pessoa jurídica} têm CEP completo no registro, geocodificável
a nível de rua. Para esse segmento --- comércio e indústria, com arranjos medianos de
66 e 92\,m$^2$ --- o chip dirigido é viável, e o custo volta para a ordem de AUD 0,20
por instalação.

\bloco{E não é onde está a análise de energia}
A análise energética por satélite que existe e é bancável usa satélite
\textbf{meteorológico geoestacionário} --- índice de nuvem $\rightarrow$ GHI
$\rightarrow$ modelo de desempenho --- a 1--2\,km e 10 minutos. Ela não olha para o
painel, olha para a nuvem. É o ramo em que as Figs.\,5 e 6 já operam.

\bloco{Ressalvas}
\ressalva{Preços e sensores são os publicados pelo fornecedor, em AUD e sem tributo; mudam sem aviso.}
\ressalva{A profundidade de arquivo não é publicada por sensor --- decisivo, porque todo desenho da série é longitudinal.}
\ressalva{Catálogo só óptico: sem SAR, que seria o complemento imune a nuvem.}
}{{\color{claro}Página auxiliar --- ver Fig.\,2 para a curva de detectabilidade.}}

%================================================================ APLICAÇÃO
\cab{APLICAÇÃO \quad · \quad exercício fora da série numerada}
{E se eu quisesse construir ao lado?}
{A série encadeada para responder a uma pergunta de expansão: desenhar uma área nova junto a Sol do Cerrado, estimar quanto ela produziria, e o que o despacho faria com essa energia.}

\begin{minipage}[t]{0.315\linewidth}\vspace{0pt}\small
\bloco{A corrente, com fatores medidos}
Nenhum dos fatores é suposto; todos saem de figuras anteriores.
\par\vspace{2pt}\scriptsize
\begin{tabular}{@{}lll@{}}
\toprule
passo & fator & origem \\
\midrule
terreno $\rightarrow$ módulo & $\div$\,2,5 & Fig.\,3 \\
módulo $\rightarrow$ capacidade & $\div$\,5,12\,m²/kWp & Fig.\,2 \\
capacidade $\rightarrow$ potencial & $\times$\,0,229 de FC & Figs.\,5 e 6 \\
\bottomrule
\end{tabular}
\par\small\vspace{2pt}
O fator de capacidade potencial vem da razão de desempenho limpa de Sol do Cerrado
(0,931, calibrada em 279 dias sem restrição) e do GHI da NASA POWER sobre Jaíba
(5,908\,kWh/m²/dia).

\bloco{O que a ida e volta prova --- e o que não prova}
Aplicada à usina que já existe, a corrente fecha: 681,3\,MW fiscalizados implicam
3,49\,km² de módulo, e o caminho inverso devolve 681,3\,MW. \textbf{Isso não é
validação}: é dividir e multiplicar pela mesma constante de 5,12\,m²/kWp, e fecharia
com qualquer valor dela. O que a ida e volta confere é a aritmética.

O elo que \textbf{não} está validado é o primeiro. A razão terreno/módulo de 2,5 é
típica de projeto, não medida nesta usina, e entra linearmente na capacidade: errá-la
em 20\% erra a capacidade em 20\%. Medi-la exigiria delinear o arranjo por
segmentação --- o mesmo passo que a Fig.\,7 deixou em aberto.

\bloco{Mas ela estima o potencial, não o entregue}
Para 2025, o método prevê fator de capacidade de \textbf{0,229} e
\textbf{1.368\,GWh}. A usina entregou fator de \textbf{0,150} e
\textbf{897\,GWh}.
\end{minipage}\hfill
\begin{minipage}[t]{0.315\linewidth}\vspace{0pt}\small
\bloco{E a lacuna é a medida}
A diferença é de \textbf{34,4\%}. O ONS apurou, por outro caminho --- curva de
desempenho declarada da usina e apuração operativa de \textit{constrained-off} ---
uma taxa de \textbf{32,8\%} na mesma usina e ano. Diferem em \textbf{1,6 ponto}.

\bloco{Quão independentes são os dois números}
Menos do que eu havia escrito, e é preciso dizer onde. A corrente \textbf{não} usa a
geração de referência do operador --- essa é a parte genuinamente externa. Mas usa a
razão de desempenho limpa de 0,931, e ela é calibrada nos dias que a
\textit{própria apuração do ONS} marca como sem restrição. O sinalizador de restrição
é compartilhado; o que se testa é a estimativa de \textit{quanto}, não a de
\textit{quando}.

E não é um terceiro número: $1 - PR_{2025}/PR_{\text{limpo}}$ é o mesmo estimador da
Fig.\,6, que dá \textbf{33,5\%} --- a diferença para os 34,4\% é só de agregação do
GHI. \textbf{São duas estimativas, não três}, e o que a página acrescenta à Fig.\,6 é
o passo área\,$\rightarrow$\,capacidade, que é justamente o não validado.

\bloco{Cenários para uma área nova}
\par\vspace{1pt}\scriptsize
\begin{tabular}{@{}rrrrr@{}}
\toprule
AOI & capac. & potencial & a 31\% & a 62\% \\
km² & MW & GWh/ano & GWh & GWh \\
\midrule
 1 &    78 &   157 &   108 &    60 \\
 2 &   156 &   314 &   216 &   119 \\
 5 &   391 &   784 &   541 &   298 \\
10 &   781 & 1.569 & 1.082 &   596 \\
20 & 1.563 & 3.137 & 2.165 & 1.192 \\
\bottomrule
\end{tabular}
\par\small\vspace{2pt}
As duas últimas colunas são o ponto: \textbf{não se pode usar a taxa média para uma
usina nova}. Novan e Wang (2024) mediram que a taxa \textbf{marginal} --- quanto de
uma usina nova é cortado --- é cerca de duas vezes a média. Daí a coluna de 62\%.
\end{minipage}\hfill
\begin{minipage}[t]{0.315\linewidth}\vspace{0pt}\small
\bloco{O que o despacho faria com essa energia}
O barramento de Jaíba já tem \textbf{1.301\,MW} em três conjuntos --- Sol do Cerrado
(681), Jaíba V (500) e Jaíba 138\,kV (120) --- todos cortados entre 24\% e 31\% na
janela de referência.

E o corte ali é \textbf{89,5\% de origem sistêmica}: 61,8\% por razão energética,
16,2\% por confiabilidade e 11,5\% por indisponibilidade externa, todos sistêmicos;
apenas 10,5\% tem origem local. Não é o barramento que está saturado --- é o sistema
que já tem oferta sobrando naquele instante. No agregado nacional a distribuição é a
mesma em espírito: 68,5\% energética, 22,0\% confiabilidade, 9,5\% indisponibilidade
externa.

Isso encadeia com as Figs.\,8 a 10: térmica inflexível nas horas de corte, corredor
\textit{vazio} e distância à demanda como preditor. Jaíba está a \textbf{605\,km de
rede} do centro de carga mais próximo --- terceiro quartil da amostra.

\bloco{Uma honestidade que os dados impõem}
Seria intuitivo dizer que mais capacidade no mesmo nó gera mais corte. \textbf{Os
dados não sustentam isso}: na Fig.\,10 a capacidade renovável no nó deu
$\rho = -0{,}24$ e não sobreviveu a Holm --- nem o grau, nem a tensão, nem o número
de usinas. O que sobreviveu foi a distância. A leitura defensável não é
``Jaíba está lotado'', é \textbf{``Jaíba está longe''} --- e uma usina nova ali
herdaria a posição, não a melhoraria.

\end{minipage}

\vspace{5pt}{\color{claro}\rule{\linewidth}{0.4pt}}\par\vspace{4pt}
\begin{minipage}[t]{0.485\linewidth}\vspace{0pt}\small
\bloco{O que pode fazer o número errar}
\ressalva{Este cálculo não usa limite de linha. Eu havia escrito que ele não seria público; o ONS publica capacidade operativa em MVA, com variantes sazonais, em 86,7\% das linhas. Incorporá-la tornaria a estimativa de exposição ao corte bem mais firme.}
\ressalva{A razão marginal de 2× é medida na Califórnia. Não há equivalente brasileiro publicado; é ordem de grandeza declarada, não um fator calibrado aqui.}
\ressalva{A razão terreno/módulo de 2,5 é o elo não validado da corrente, discutido acima; inclinação, rastreamento e relevo a deslocam.}
\ressalva{O fator de capacidade potencial usa o desempenho limpo de Sol do Cerrado; uma usina nova com outro módulo ou rastreamento teria outro.}
\end{minipage}\hfill
\begin{minipage}[t]{0.485\linewidth}\vspace{0pt}\small
\bloco{Escopo}
O que sai daqui é uma estimativa \textbf{física e operativa} --- quanto a área
produziria e quanto o despacho absorveria. O que ela decide depende do uso: para
dimensionar, comparar sítios ou antecipar exposição ao corte, basta. Para uma decisão
de investimento entram contrato, outorga e parecer de acesso --- outra camada, que
não está em dado aberto e que este exercício não tenta cobrir.
\par\vspace{3pt}
{\color{claro}\scriptsize Dados-fonte em \texttt{data/processed/aplicacao\_*.csv},
gerados por \texttt{notebooks/\_aplicacao\_expansao.py}.}
\end{minipage}
\clearpage

%================================================================ FUNDAMENTAÇÃO
\cab{FUNDAMENTAÇÃO}{Onde esta série se encaixa}{O que já estava publicado, o que a série confirma, o que ela acrescenta --- e o que mudou na regulação enquanto a análise era feita.}

\begin{minipage}[t]{0.315\linewidth}\vspace{0pt}\footnotesize
\bloco{O Brasil está no topo da faixa internacional}
\scriptsize
\begin{tabular}{@{}lll@{}}
\toprule
sistema & taxa & referência \\
\midrule
Brasil, solar (apurada) & \textbf{21,6\%} & \textit{esta série} \\
Brasil, eól.+solar 2025 & 20,7\% & ONS \\
Chile 2024 & 17\% & imprensa setorial \\
Brasil, eólica (apurada) & \textbf{15,9\%} & \textit{esta série} \\
Irlanda, eólica 2025 & 11,3\% & EirGrid \\
Espanha, jul/2025 & $\approx$11\% & recorde \\
Califórnia, solar & 4,3\% & Novan e Wang \\
\bottomrule
\end{tabular}
\par\vspace{3pt}\footnotesize
Os denominadores diferem e a comparação com os sistemas estrangeiros é de ordem de
grandeza, não pareada. As \textbf{duas linhas brasileiras, essas sim, são pareadas}:
ambas são razão de energia --- cortada dividida por disponível --- na janela
abr/2024--ago/2026, em que solar e eólica têm apuração simultânea, sobre todos os
conjuntos com apuração (87 solares, 178 eólicos). Numa única extração de hoje elas dão
22,1\% e 15,9\%; a linha da solar preserva os 21,6\% da Fig.\,1, apurados em extração
anterior.

O resultado é claro e vale para as duas fontes: o corte brasileiro está no topo da
faixa internacional, com a solar acima do Chile e a \textbf{eólica acima da Irlanda}.
A diferença entre as duas fontes existe --- a solar perde cerca de 1,4 vez o que a
eólica perde --- mas é bem menor do que sugeriria comparar a taxa agregada de uma com
a mediana da outra. O que separa os dois regimes não é o \textit{tamanho} do corte, e
sim a \textbf{composição da causa}: 91\% de origem sistêmica na solar contra 76\% na
eólica. É esse o achado da Fig.\,11.

E há um agravante de leitura. Novan e Wang (2024) mostram que a taxa
\textbf{marginal} --- quanto de uma usina \textit{nova} será cortado --- é cerca de
duas vezes a média. Na Califórnia, média de 4,3\% contra marginal de 9\%. Se a razão
se mantiver, a exposição de um projeto brasileiro entrando agora é muito superior
aos 21,6\% medidos.

\bloco{O mecanismo não é inédito --- o teste é}
A competição entre mínimo técnico térmico e renovável variável é o problema central
da literatura chinesa de curtailment, e gerou uma resposta de política explícita: o
programa de \textit{flexibility retrofit}, com 16,35\,GW em demonstração em 2017.

O que não existia era o teste com a decomposição de despacho publicada pelo ONS,
que separa geração por ordem de mérito de geração por inflexibilidade declarada,
razão elétrica, \textit{unit commitment} e segurança energética. As Figs.\,8 e 9
usam esse campo --- e a Fig.\,9 mostra que ele não sustenta atribuição local.
\end{minipage}\hfill
\begin{minipage}[t]{0.315\linewidth}\vspace{0pt}\footnotesize
\bloco{O que a série confirma}
\textbf{MMGD como motor do corte.} Vieira et al.\ (2026) acharam a MMGD como
preditor mais forte do corte por demanda ($OR\approx9{,}8$ para solar), com a carga
do sistema atuando de forma protetiva. O gradiente por quintil de carga da Fig.\,8
reproduz isso de forma independente.

\textbf{Invisibilidade da geração atrás do medidor.} É problema documentado: a
geração distribuída não é observada pelo operador e eleva o erro de previsão de
carga líquida em 30\% no nível domiciliar e até 250\% no agregado. A Fig.\,2 mede a
outra face do mesmo problema --- ela também não é observável por satélite gratuito.

\textbf{Resolução necessária para telhado.} Toda a literatura que detecta MMGD com
sucesso opera entre 0,1 e 0,3\,m. O inventário global de Kruitwagen et al.\ (2021),
o trabalho de referência com Sentinel-2, cobre apenas instalações acima de 10\,kW.

\bloco{O que ela acrescenta}
\textbf{Uma validação externa da referência do ONS.} A geração de referência é
derivada de curvas de desempenho da própria usina. Validá-la exige fonte
independente --- feito no lado eólico com anemometria de SCADA (Silva et al., 2026)
e, aqui, no lado solar com irradiância de satélite, em 74 usinas.

\textbf{Um resultado negativo sobre observabilidade.} A obra é detectável e datável
por Sentinel-2; a operação não é. Isso inverte a intuição de auditoria por imagem.

\textbf{A refutação de uma leitura própria.} A Fig.\,9 mostra que a associação entre
corte e térmica inflexível não distingue causa local de fator comum --- e retira da
Fig.\,8 uma atribuição causal que eu havia feito.

\textbf{E o candidato a contribuição original.} A literatura documenta que as duas
causas existem (Bird et al.), mede a proporção entre elas (EirGrid) e argumenta que o
corte deve ser estimado por nó (Maji et al.). A Fig.\,11 dá o passo seguinte: mostra
que \textbf{a covariável que prevê o corte muda conforme a composição da causa} ---
preditor sistêmico onde o corte é sistêmico, preditor local onde é local. Não
encontrei isso publicado; ``não encontrei'' não é ``não existe'', e a busca não foi
sistemática.
\end{minipage}\hfill
\begin{minipage}[t]{0.315\linewidth}\vspace{0pt}\footnotesize
\bloco{A regulação se moveu durante a análise}
A questão que a Fig.\,3 deixou explicitamente em aberto --- se o corte é ressarcido
--- deixou de ser hipotética.

A \textbf{Lei 15.269/2025} criou compensação por \textit{constrained-off} para
eólicas e solares, custeada por Encargos de Serviços do Sistema, com efeitos
retroativos a 1\textsuperscript{o} de setembro de 2023 e condicionada à renúncia de
ações judiciais e à assinatura de termo de compromisso. A Portaria MME 140
regulamentou os cortes ocorridos até 25 de novembro de 2025; os posteriores ficam
com a ANEEL.

A primeira fase do processo reuniu manifestação de interesse de \textbf{1.539
usinas}, cerca de 53\,GW instalados. Estimativas setoriais apontam perda da ordem de
\textbf{R\$\,6,5 bilhões} em 2025 com \textit{constrained-off} de eólica e solar ---
mesma ordem de grandeza da conta a CMO da Fig.\,3, extrapolada para o parque.

Isso reforça a ressalva daquela figura: CMO mede valor sistêmico, não ressarcimento.
Com a lei, os dois passam a divergir por construção --- o gerador é compensado por
energia cujo valor sistêmico, no instante do corte, era próximo de zero.

\bloco{Nota sobre a qualidade das fontes}
Nem tudo aqui tem o mesmo peso. São revisadas por pares: Vieira et al.\ (2026),
Novan e Wang (2024), Kruitwagen et al.\ (2021), Yu et al.\ (2018), Jiang et al.\
(2021), Alix-García e Millimet (2023). São \textit{working papers} ou pré-prints:
Newbery (2024), Proctor et al.\ (2023), Silva et al.\ (2026), Global Renewables
Watch. São fontes normativas: a Lei 15.269/2025 e a Portaria MME 140. E são imprensa
setorial as taxas de Chile, Espanha e Europa e o valor de R\$\,6,5\,bi --- usadas
como ordem de grandeza, nunca como evidência.

\end{minipage}
\clearpage

%================================================================ REFERÊNCIAS
\cab{REFERÊNCIAS}{Fontes citadas}{Bibliografia, bases de dados e atos normativos.}

\begin{minipage}[t]{0.485\linewidth}\vspace{0pt}\footnotesize
\bloco{Curtailment e despacho}
Vieira, G.T.T.; Silva, W.N.; Lourenço, L.F.N.; Monaro, R.M.; Salles, M.B.C.;
Almeida, C.F.M. \textit{Characterizing wind and solar curtailment in Brazil: an
evidence-based analysis of operational drivers.} Electric Power Systems Research,
2026. \par\vspace{2pt}
Novan, K.; Wang, Y. \textit{Estimates of the marginal curtailment rates for solar and
wind generation.} Journal of Environmental Economics and Management, 2024.
\par\vspace{2pt}
Newbery, D. \textit{Marginal curtailment of wind and solar PV: transmission
constraints, pricing and access regimes for efficient investment.} EPRG Working
Paper 2401, University of Cambridge, 2024. \par\vspace{2pt}
Silva, W.N. et al. \textit{An integrated methodology for assessing wind power
curtailment using anemometric measurements and operational data in the Brazilian
context.} Atmosphere, 2026. \par\vspace{2pt}
Bird, L.; Lew, D.; Milligan, M. et al. \textit{Wind and solar energy curtailment: a
review of international experience.} Renewable and Sustainable Energy Reviews, 65,
577--586, 2016. \par\vspace{2pt}
Maji, D.; Irwin, D.E.; Shenoy, P.J.; Sitaraman, R.K. \textit{A first look at
node-level curtailment of renewable energy and its implications.} ACM e-Energy '25,
Rotterdam, 2025. \par\vspace{2pt}
EirGrid. \textit{Annual renewable energy constraint and curtailment report}, edições
2013--2024. \par\vspace{2pt}
Clean Energy Ministerial. \textit{Thermal power plant flexibility.} Advanced Power
Plant Flexibility Campaign, 2018.

\bloco{Sensoriamento remoto de fotovoltaica}
Kruitwagen, L.; Story, K.T.; Friedrich, J.; Byers, L.; Skillman, S.; Hepburn, C.
\textit{A global inventory of photovoltaic solar energy generating units.} Nature,
598, 604--610, 2021. \par\vspace{2pt}
Yu, J.; Wang, Z.; Majumdar, A.; Rajagopal, R. \textit{DeepSolar: a machine learning
framework to efficiently construct a solar deployment database in the United
States.} Joule, 2(12), 2018. \par\vspace{2pt}
Jiang, H. et al. \textit{Multi-resolution dataset for photovoltaic panel segmentation
from satellite and aerial imagery.} Earth System Science Data, 13, 5389--5401, 2021.
\par\vspace{2pt}
Robinson, C. et al. \textit{Temporal cluster matching for change detection of
structures from satellite imagery.} Microsoft Research, 2021. \par\vspace{2pt}
Microsoft Research. \textit{Global Renewables Watch: a temporal dataset of solar and
wind energy derived from satellite imagery.} arXiv:2503.14860, 2025.
\end{minipage}\hfill
\begin{minipage}[t]{0.485\linewidth}\vspace{0pt}\footnotesize
\bloco{Método --- variáveis derivadas de sensoriamento}
Proctor, J.; Carleton, T.; Sum, S. \textit{Parameter recovery using remotely sensed
variables.} NBER Working Paper 30861, 2023. \par\vspace{2pt}
Alix-García, J.; Millimet, D.L. \textit{Remotely incorrect? Accounting for
nonclassical measurement error in satellite data on deforestation.} Journal of the
Association of Environmental and Resource Economists, 10(5), 1335--1367, 2023.
\par\vspace{2pt}
Callaway, B.; Sant'Anna, P.H.C. \textit{Difference-in-differences with multiple time
periods.} Journal of Econometrics, 225(2), 2021.

\bloco{Atos normativos}
Lei n\textsuperscript{o} 15.269, de 2025 --- compensação por \textit{constrained-off}
de fontes eólica e solar. \par\vspace{2pt}
Portaria MME n\textsuperscript{o} 140 --- procedimentos de compensação para cortes
entre 1/9/2023 e 25/11/2025. \par\vspace{2pt}
Lei n\textsuperscript{o} 14.300, de 2022 --- marco legal da microgeração e
minigeração distribuída.

\bloco{Bases de dados}
EPE, webmap institucional (serviço ArcGIS REST): camadas de linhas de transmissão e
subestações em operação, com geometria. Alternativa tabular: ANEEL, SIGET ---
\textit{Contrato Módulo Linha Transmissão Subestação Origem Subestação Destino}.
\par\vspace{2pt}
ONS, Portal de Dados Abertos (\texttt{dados.ons.org.br}): restrição por
\textit{constrained-off} fotovoltaico e eólico; geração térmica por motivo de
despacho; fator de capacidade; curva de carga; intercâmbio nacional; energia
armazenada por subsistema; CMO semi-horário; relação conjunto--usina.
\par\vspace{2pt}
ANEEL: Sistema de Informações de Geração (SIGA); relação de empreendimentos de
geração distribuída e informações técnicas fotovoltaicas. \par\vspace{2pt}
ESA/Copernicus: Sentinel-2 L2A, via catálogo STAC público da AWS
(\texttt{earth-search.aws.element84.com}). \par\vspace{2pt}
NASA POWER: irradiância diária (MERRA-2 e SYN1DEG). \par\vspace{2pt}
IBGE: malhas territoriais das unidades da federação.

\bloco{Fontes de imprensa setorial}
Taxas de \textit{curtailment} de Chile, Espanha e Europa e estimativa de perda
financeira no Brasil: \textit{pv magazine}, Rystad Energy, Canal Solar, EIA
\textit{Today in Energy}. Usadas como ordem de grandeza.
\end{minipage}
\clearpage

%================================================================ SÍNTESE
\cab{FECHAMENTO}{O que ficou de pé}{Estabelecido, derrubado e em aberto, ao fim de doze análises.}

\begin{minipage}[t]{0.315\linewidth}\vspace{0pt}\small
\bloco{Estabelecido}
\textbf{O corte é grande e sistêmico.} 21,6\% no parque com apuração, dois terços
por razão energética, incidindo sobre toda a janela solar.

\textbf{A referência do ONS é bem calibrada.} Viés mediano 1,02 em 74 usinas, sem
piorar onde o corte é mais intenso. O agregado da Fig.\,1 é confiável.

\textbf{A perda de Sol do Cerrado é de rede, não de clima.} Nove décimos da queda de
29\% no fator de capacidade não se explicam pela irradiância.

\textbf{A MMGD é sub-pixel.} 91\% abaixo de um pixel Sentinel-2; verificação por
imagem exige 2\,m ou melhor, e coordenada por instalação, que não existe.

\textbf{A obra é detectável e datável; a operação não.} +34\% de brilho na
construção ($p=1{,}2\times10^{-5}$) contra nada na operação ($p=0{,}09$).

\textbf{O piso térmico inflexível é real e está no Sudeste.} 1.739\,MWmed de
inflexibilidade pura, contra 36 no Nordeste.

\textbf{O grau do nó prevê o corte eólico.} $\rho = +0{,}35$, único dos quatro achados
de rede a passar em bootstrap, partição estrutural e troca de grafo (Fig.\,12). É o
resultado mais firme da série sobre alocação de corte.

\textbf{E o preditor espelha a origem do corte.} Solar é 91\% sistêmico e responde à
distância à demanda; eólica tem 24\% de corte local e responde a aglomeração no nó
($\rho = +0{,}47$). A dissociação sobrevive ao controle de geografia.

\textbf{O corte pode ser medido sem a referência do operador.} Irradiância mais razão de
desempenho calibrada estimam 33,5\% de perda em Sol do Cerrado em 2025 (Fig.\,6), e
34,4\% com o passo de área acoplado, contra 32,8\% apurados pelo ONS --- 1,6 ponto de
diferença. Independência \textbf{parcial}: dispensa a geração de referência, mas ainda
usa a apuração do ONS para escolher os dias de calibração.
\end{minipage}\hfill
\begin{minipage}[t]{0.315\linewidth}\vspace{0pt}\small
\bloco{Derrubado --- inclusive por mim}
\textbf{``A referência do ONS infla o corte''} (Fig.\,5). Era artefato do
contrafactual: 2023 já era um ano restrito. Corrigido na Fig.\,6, que \textit{reforça}
a conclusão principal --- corte implícito de 33\%, não 26\%.

\textbf{``A térmica inflexível disputa espaço com a solar''} (Fig.\,8). A
decomposição da Fig.\,9 mostra que essa térmica se associa igualmente ao corte a
dois mil quilômetros: é marcador sistêmico, e o desenho não separa as hipóteses.

\textbf{``O Nordeste é cortado porque o corredor satura''.} É o inverso: é cortado
com o fluxo NE$\rightarrow$SE vazio.

\textbf{``A térmica não abre espaço para a solar''.} Não sobrevive ao controle de
carga --- a carga líquida diária é plana.

\textbf{``Explicações locais do corte generalizam''} (Fig.\,3). Nenhuma das cinco
covariáveis testadas se reproduz no parque.

\textbf{``A distância à demanda é sobre posição na rede, não sobre tecnologia''}
(Fig.\,10). A Fig.\,11 mostra que não: vale para fonte de corte sistêmico. E a
Fig.\,12 mostra que, mesmo ali, é frágil --- cai de 0,62 para 0,11 conforme a metade
da amostra, e some no grafo por barra.

\textbf{``Impedância e limite de linha não estão em dado aberto''} (Figs.\,3, 10, 11 e
aplicação). Estão, no pacote \texttt{linha-transmissao} do ONS: reatância em 100\% das
linhas e capacidade em MVA em 87\%. Afirmei o contrário quatro vezes sem verificar.

\textbf{``Detectar painéis por satélite audita política de MMGD''} --- a premissa
original. Fechada pela Fig.\,2.
\end{minipage}\hfill
\begin{minipage}[t]{0.315\linewidth}\vspace{0pt}\small
\bloco{Em aberto}
\textbf{Por que \textit{esta} usina e não aquela? --- parcialmente respondido.} A
Fig.\,10 achou duas covariáveis que sobrevivem à correção de multiplicidade:
distância de rede à demanda e capacidade da usina. Sob os testes da Fig.\,12 é a
\textbf{segunda} que resiste --- e é a que não tem explicação testada. A variância
explicada é modesta, e a maior parte da dispersão da taxa de corte segue sem
explicação.

\textbf{O que causa o corte do Nordeste.} Nem piso térmico local, nem escoamento.
Resta ausência de demanda no destino, por eliminação.

\textbf{A inflexibilidade declarada é o mínimo técnico real?} É a quantidade que uma
auditoria precisaria, e não está publicada.

\bloco{Próximos passos, em ordem de retorno}
\textbf{1. Reconstruir o corte anterior a abr/2024.} A maquinaria da Fig.\,6 já
existe; um modelo treinado nos intervalos sem restrição, aplicado para trás,
destrava o teste temporal da Fig.\,4. Há verdade-terreno para validar.

\textbf{2. Medir o resíduo inexplicável.} No nível de intervalo, a diferença de
desempenho entre um modelo só com observáveis sistêmicos e outro com identidade de
usina quantifica o que as Figs.\,10 e 11 deixaram por explicar. A série eólica
acrescenta 1,4\,GB e três anos a mais de painel.

\textbf{3. Trocar NASA POWER por NSRDB.} De célula de $\sim$50\,km e passo diário
para 2\,km e 10 minutos, gratuito. Resolve três ressalvas já escritas.

\textbf{4. Delinear o arranjo por segmentação.} Substitui os círculos de 400\,m da
Fig.\,7 e aperta a datação.

\bloco{Uma regra que vale para 1 e 2}
Qualquer saída de modelo que entre numa inferência a jusante enviesa ponto e erro
padrão. Predição entra como distribuição, nunca como valor pontual.
\end{minipage}

\vfill
{\color{claro}\rule{\linewidth}{0.4pt}}\par\vspace{3pt}
{\color{claro}\scriptsize
Todas as análises usam exclusivamente dados abertos e sem autenticação. Cada figura
exporta seus dados-fonte em CSV e passa por auditoria de tamanho mínimo de glifo e de
colisão de elementos antes de ser considerada pronta. Notebooks e scripts geradores
estão em \texttt{notebooks/}; figuras em \texttt{../figures/}; dados-fonte em
\texttt{data/processed/}.}

\end{document}
